\lysh{33662}{4}
\begin{alist}
    \item Η συνάρτηση $f$ είναι πολυωνυμική, άρα παραγωγίσιμη στο διάστημα $(0, +\infty)$, με παράγωγο:
    \[ f'(x) = 4x^3 - 4x, x > 0. \]
    Ισχύει ότι $f'(x) = 4x(x^2 - 1) = 4x(x-1)(x+1)$ και $f'(x) = 0 \Leftrightarrow 4x(x^2 - 1) = 0 \Leftrightarrow x = 0$ ή $x = -1$ ή $x = 1$, από τις οποίες γίνεται δεκτή μόνο η $x = 1$, εφόσον $x > 0$.
    
    Επίσης, εφόσον $x > 0$, $f'(x) < 0 \Leftrightarrow 4x(x^2 - 1) < 0 \Leftrightarrow x^2 - 1 < 0 \Leftrightarrow 0 < x < 1$ και
    $f'(x) > 0 \Leftrightarrow 4x(x^2 - 1) > 0 \Leftrightarrow x^2 - 1 > 0 \Leftrightarrow x > 1$.
    
    Οπότε, ισχύει ότι $f'(x) < 0 \Leftrightarrow x \in (0,1)$ και $f'(x) > 0 \Leftrightarrow x > 1$.
    
    \item Σύμφωνα με το ερώτημα (α), η συνάρτηση $f$ είναι γνησίως φθίνουσα στο $(0,1]$ και γνησίως αύξουσα στο $[1, +\infty)$, οπότε έχει ελάχιστη τιμή στο $1$, το $f(1) = 1$.
    
    \item Η συνάρτηση $f$ έχει ελάχιστη τιμή το $1$, οπότε θα ισχύει ότι $f(x) \geq 1 > 0 \Rightarrow f(x) > 0$, για κάθε $x > 0$.
\end{alist}
\vspace{6mm}
\noindent\rule{\textwidth}{0.4pt}
\vspace{6mm}

